% !TeX program = xelatex
% Nakupni prehled v sazbe TUL 2.3.
\PassOptionsToPackage{localfonts}{tul}
\documentclass[FM,SP,oneside]{tulthesis}
\usepackage{polyglossia}
\setdefaultlanguage{czech}
\setotherlanguage{english}
\usepackage{csquotes}
\usepackage{microtype}
\usepackage{booktabs}
\usepackage{longtable}
\usepackage{array}
\usepackage{tabularx}
\usepackage{siunitx}
\usepackage{enumitem}
\usepackage{xcolor}
\usepackage{tikz}
\usepackage{xurl}
\usetikzlibrary{arrows.meta,positioning,shapes.geometric,fit}
\usepackage[backend=biber,style=numeric,sorting=none,maxbibnames=20]{biblatex}
\addbibresource{references.bib}
\AtBeginBibliography{\small}

% Okraje pro tisk a vazbu: vnitrni 30 mm, vnejsi 20 mm; sazba A4 siroka 160 mm.
% V rezimu oneside je vazba vlevo na kazdem listu. Rezim twoside zrcadli
% okraje sudych stran. Horni/dolni rozmery ponechavame podle TUL 2.3.
% LaTeX meri bocni posun od 1in = 25,4 mm.
\setlength{\oddsidemargin}{4.6mm}
\makeatletter
\if@twoside
  \setlength{\evensidemargin}{-5.4mm}
\else
  \setlength{\evensidemargin}{4.6mm}
\fi
\makeatother
\setlength{\textwidth}{160mm}
\setlength{\emergencystretch}{2em}
% Neprotahovat svisle mezery pri prepnuti reportu do oboustranne sazby.
\raggedbottom

% Uprednostni zalomeni dlouhych URL v bibliografii pred pretekanim do okraje.
\setcounter{biburlnumpenalty}{9000}
\setcounter{biburlucpenalty}{9000}
\setcounter{biburllcpenalty}{9000}

\sisetup{locale = DE, per-mode = symbol, range-phrase = { až }}
\DeclareSIUnit\byte{B}
\setlist{itemsep=0.25em,topsep=0.45em}
\newcolumntype{Y}{>{\raggedright\arraybackslash}X}
\newcolumntype{L}[1]{>{\raggedright\arraybackslash}p{#1}}
\newcolumntype{R}[1]{>{\raggedleft\arraybackslash}p{#1}}
\newcommand{\todoauthor}[1]{\textcolor{red}{[DOPLNIT: #1]}}


\usepackage{colortbl}
\usepackage{ragged2e}
\hypersetup{
  colorlinks=true,
  urlcolor=tulFM,
  linkcolor=tulFM,
  pdfauthor={Jiří Buršík},
  pdftitle={Nákupní seznam - RTK GNSS/IMU systém}
}

\setlength{\parindent}{0pt}
\setlength{\parskip}{0pt}
\renewcommand{\arraystretch}{1.08}
\pagestyle{fancy}
\fancyhf{}
\renewcommand{\headrulewidth}{0pt}
\fancyfoot[L]{\color{tulFM}\footnotesize Technická univerzita v Liberci -- FM}
\fancyfoot[R]{\color{tulFM}\footnotesize\thepage/2}

\newcommand{\cartheader}{%
  \noindent
  \begin{minipage}[c]{0.76\textwidth}
    \includegraphics[width=52mm]{logo-fm-txt.pdf}
  \end{minipage}%
  \hfill
  \begin{minipage}[c]{0.18\textwidth}
    \raggedleft\includegraphics[width=14mm]{symbol-fm.pdf}
  \end{minipage}
  \vspace{1mm}
  {\color{tulFM}\rule{\textwidth}{0.9pt}}
  \vspace{1mm}
}

\newcommand{\shopheading}[2]{%
  \vspace{0.55em}
  \noindent\colorbox{tulFM!10}{%
    \parbox{\dimexpr\textwidth-2\fboxsep\relax}{%
      \textcolor{tulFM}{\TULmono\bfseries\normalsize #1}%
      \hfill
      {\small\bfseries #2}%
    }%
  }
  \vspace{0.2em}
}

\newcommand{\product}[3]{%
  \href{#3}{\textbf{#1}}\par\vspace{-0.25em}%
  {\scriptsize\color{black!62}#2\par}
}

\begin{document}
\cartheader

{\sffamily\bfseries\fontsize{16pt}{19pt}\selectfont\color{tulFM}
Nákupní seznam -- RTK GNSS/IMU systém\par}
\vspace{0.2em}
{\footnotesize
\textbf{Projekt:} měření polohy vozidla Formula Student\hfill
\textbf{Autor:} Jiří Buršík\\
\textbf{Aktualizace seznamu:} 30. 9. 2026\hfill
\textbf{Cílový rozpočet:} 25\,000 Kč bez DPH
}

\shopheading{RPiShop.cz}{3\,932 Kč s DPH}

\footnotesize
\begin{tabularx}{\textwidth}{@{}>{\centering\arraybackslash}p{10mm}Y>{\raggedleft\arraybackslash}p{29mm}@{}}
\rowcolor{tulFM!7}
\textbf{Ks} & \textbf{Položka} & \textbf{Cena s DPH} \\
\midrule
1 & \product{Raspberry Pi 5 -- 2 GB RAM}{Kód RPI502; výpočetní jednotka roveru.}{https://rpishop.cz/535838/raspberry-pi-5-2gb-ram/} & 1\,629 Kč \\
1 & \product{Raspberry Pi 5 Active Cooler}{Kód OFI118; aktivní chlazení pro provoz v krytu.}{https://rpishop.cz/chladice-pro-raspberry-pi-5/6496-raspberry-pi-5-active-cooler.html} & 149 Kč \\
1 & \product{Raspberry Pi M.2 HAT+}{Kód OFI134; pro vlastní NVMe SSD formátu 2230/2242.}{https://rpishop.cz/502557/raspberry-pi-m-2-hat/} & 349 Kč \\
4 & \product{XIAO ESP32S3 + Wio-SX1262 Kit}{Kód SEE319; sada Seeed Studio pro LoRa 862--930 MHz, včetně LoRa a Wi-Fi antény.}{https://rpishop.cz/595653/xiao-esp32s3-wio-sx1262-kit-pro-meshtastic-lora/} & 4 $\times$ 339 Kč \\
1 & \product{Adafruit ICM-20948, 9osá IMU}{Kód ADA076; akcelerometr, gyroskop a magnetometr. SPI/I\textsuperscript{2}C, vývod INT, STEMMA QT/Qwiic.}{https://rpishop.cz/652161/adafruit-9osy-imu-modul-icm-20948-se-stemma-qt-qwiic/} & 449 Kč \\
\bottomrule
\end{tabularx}

\shopheading{ArduSimple}{863 EUR bez DPH}

\begin{tabularx}{\textwidth}{@{}>{\centering\arraybackslash}p{10mm}Y>{\raggedleft\arraybackslash}p{29mm}@{}}
\rowcolor{tulFM!7}
\textbf{Ks} & \textbf{Položka} & \textbf{Cena} \\
\midrule
1 & \product{simpleRTK3B Pro}{Rover s přijímačem Septentrio mosaic-X5; až 100 Hz, SBF/NMEA/RTCM.}{https://www.ardusimple.com/product/simplertk3b-x5/} & 569 EUR \\
2 & \product{u-blox ANN-MB2-00}{Třípásmová L1/L2/L5 anténa; 1 kus pro rover a 1 pro základnu. Každá s 5m kabelem a SMA konektorem.}{https://www.ardusimple.com/product/u-blox-gnss-tripleband-antenna-ann-mb2-00/} & 2 $\times$ 61 EUR \\
1 & \product{simpleRTK2B Budget}{Přenosná základna s u-blox ZED-F9P.}{https://www.ardusimple.com/product/simplertk2b/} & 172 EUR \\
\bottomrule
\end{tabularx}

\vspace{0.4em}
{\scriptsize\color{black!62}
Přepočet v souhrnu používá kurz ČNB 24,380 Kč/EUR platný 23. 9. 2026.
Ceny LoRa, IMU a antény ANN-MB2-00 ověřeny 30. 9. 2026; ostatní ceny převzaty ze seznamu z 24. 9. 2026.
Kliknutím na název položky se otevře příslušná produktová stránka.
}

\newpage
\cartheader

{\sffamily\bfseries\fontsize{16pt}{19pt}\selectfont\color{tulFM}
Nákupní seznam -- souhrn a zapojení\par}

\shopheading{Souhrn nákupu}{24\,290 Kč bez DPH}

\footnotesize
\begin{tabularx}{\textwidth}{@{}Y>{\raggedleft\arraybackslash}p{38mm}@{}}
\rowcolor{tulFM!7}
\textbf{Dodavatel} & \textbf{Mezisoučet bez DPH} \\
\midrule
RPiShop.cz & 3\,249,59 Kč \\
ArduSimple (863 EUR $\times$ 24,380 Kč/EUR) & 21\,039,94 Kč \\
\midrule
\textbf{Elektronika celkem (zaokrouhleno)} & \textbf{24\,290 Kč} \\
Cílový rozpočet & 25\,000 Kč \\
\textbf{Rezerva před dopravou a instalací} & \textbf{710 Kč} \\
\bottomrule
\end{tabularx}

\vspace{0.45em}
{\scriptsize\color{black!62}
Součet nezahrnuje dopravu, napájení, USB a signálovou kabeláž, kryty, držáky ani stativ.
SSD dodá autor. Rezerva proto není konečnou úsporou celé funkční instalace.\par
}
\vspace{0.35em}
{\sffamily\bfseries\color{tulFM}Kontroly před objednáním}
\begin{itemize}[leftmargin=1.35em,itemsep=0.05em,topsep=0.2em,parsep=0pt]
  \item IMU připojím přímo k Raspberry Pi přes \textbf{SPI a INT}, desku napájím 3,3 V. STEMMA QT/Qwiic je rozhraní I\textsuperscript{2}C; pro SPI použiji pájecí vývody.
  \item Pro akcelerometr a gyroskop počítám s nominální frekvencí \textbf{562,5 Hz} při zapnutém digitálním filtru a registru děliče nastaveném na 1. Magnetometr poskytuje nejvýše \textbf{100 Hz}. Sběr přes FIFO a přerušení ověřím na Raspberry Pi OS.
  \item Vlastní NVMe SSD musí mít formát \textbf{2230 nebo 2242}; oficiální M.2 HAT+ nepodporuje 2280.
  \item XIAO připojím k Raspberry Pi přes \textbf{USB-C datový kabel}; sada kabel neobsahuje. Pro přenos RTCM počítám s vlastním firmwarem USB/UART--LoRa namísto předinstalovaného Meshtastic.
  \item Před závodem ověřím propustnost RTCM, dosah a obnovu spojení při zvoleném nastavení rádiového přenosu v pásmu 868 MHz.
  \item Napájecí měnič se vybere až podle rozsahu palubního napětí, proudové rezervy a požadované ochrany proti přepětí a přepólování.
\end{itemize}

\vspace{0.4em}
{\sffamily\bfseries\color{tulFM}Volba IMU pro prototyp}
\par\vspace{0.2em}
ICM-20948 volím jako cenově dostupnou 9osou variantu od stejného dodavatele.
Akcelerometr umožňuje rozsah až $\pm16\,g$, gyroskop až $\pm2000\,{}^{\circ}/\mathrm{s}$.
Počítám s vyšší katalogovou hustotou šumu gyroskopu než u původního ISM330DHCX
(typicky 0,015 oproti 0,005 $({}^{\circ}/\mathrm{s})/\sqrt{\mathrm{Hz}}$ za podmínek příslušných datasheetů).
Magnetometr budu používat až po kalibraci a ověření rušení ve vozidle.
TDK označuje ICM-20948 jako \textbf{EOL}; pro případnou další výrobu bude nutné řešit náhradu.

\vspace{0.3em}
{\scriptsize\color{black!62}
Technické podklady:
\href{https://learn.adafruit.com/adafruit-tdk-invensense-icm-20948-9-dof-imu/pinouts}{zapojení Adafruit},
\href{https://product.tdk.com/system/files/dam/doc/product/sensor/mortion-inertial/imu/data_sheet/ds-000189-icm-20948-v1.5.pdf}{datasheet ICM-20948},
\href{https://www.st.com/resource/en/datasheet/ism330dhcx.pdf}{datasheet ISM330DHCX},
\href{https://www.invensense.tdk.com/en-us/products/9-axis/icm-20948}{stav produktu TDK}.
}

\vspace{0.35em}
{\scriptsize\color{black!62}
Podklad: technická rešerše RTK GNSS/IMU systému pro Formula Student, Jiří Buršík, 2026.
}

\end{document}
